\chapter{A Spillable Slice Store for NebulaStream\label{cha:solution}}
% Drafted per POINT-SEQUENCE.md Ch.4 (points 5.1–5.11 design, 6.1–6.6 implementation).

This chapter presents the backend.
It works top-down.
The final solution comes first, at a high level.
The chapter then decomposes it into components and describes each one.
The design reached that solution over four iterations, and the later iterations exist because earlier ones fell short.
The chapter presents each iteration as an alternative, with the reason it was needed, so the reader sees the development path as the justification for the final design.
The last sections describe the implementation: the integration points, the instrumentation the system exposes, and the harness that runs the evaluation.

\section{The Solution at a Glance\label{sec:solution-overview}}

The backend stores tumbling-window state that is larger than main memory.
It replaces NebulaStream's in-memory slice store with a spillable slice store behind the same interface, so operators above it do not change (Section~\ref{sec:architecture}).
Figure~\ref{fig:system-overview} shows the components involved and the path window state takes through them; the rest of this chapter describes each component in turn.
A background spiller thread moves completed windows to a RocksDB backend on local disk before memory pressure becomes fatal.

% Conceptual system overview: the path window state takes through the components.
% Companion to figures/spillable_slice_store.tex (class-level integration view):
% this figure is the DATA-FLOW view for the non-expert reader, placed in Section 4.1.
% Same color legend as fig:ooc-integration: grey = existing NES, blue = this thesis,
% orange = storage engine / persistent storage.
% Preamble requirements already in Template.tex (tikz + positioning,arrows.meta,fit,backgrounds,calc).

\begin{figure}[ht]
\centering
\resizebox{\textwidth}{!}{%
\begin{tikzpicture}[
  font=\sffamily\small,
  every node/.style={align=center},
  box/.style     ={rectangle, draw, thick, minimum height=1.5cm, inner sep=5pt, rounded corners=1pt, minimum width=3.2cm},
  existing/.style={box, fill=black!12},
  newcomp/.style ={box, fill=blue!18},
  storage/.style ={box, fill=orange!28},
  slice/.style   ={rectangle, draw, thick, minimum width=2.6cm, minimum height=1.5cm, inner sep=4pt, rounded corners=1pt, fill=blue!18},
  flow/.style    ={-{Latex[length=2.6mm]}, very thick},
  watch/.style   ={-{Latex[length=2.4mm]}, thick, dashed, black!60},
  lbl/.style     ={fill=white, draw=black!30, line width=0.2pt, rounded corners=1pt,
                   inner sep=2.5pt, font=\sffamily\footnotesize\itshape, text=black!75},
]

% ================= NODES =================
% ---- outside the cap, left: the input ----
\node[existing, minimum width=2.9cm] (src) at (-0.4,0)
  {Sensor stream\\\footnotesize\itshape unbounded tuples};

% ---- row A: processing path ----
\node[existing] (workers) at (4.2,0)
  {Worker threads\\\footnotesize\itshape execute pipeline tasks};
\node[slice] (hot) at (9.0,0)
  {Hot slice\\\footnotesize\itshape window still filling\\\footnotesize\itshape (never spillable)};
\node[slice, minimum width=3.4cm] (cold) at (13.8,0)
  {Cold slices ($\sim N_w$)\\\footnotesize\itshape complete, held for\\\footnotesize\itshape emit lag $L$};

% ---- outside the cap, right: results ----
\node[existing, minimum width=2.9cm] (emit) at (19.6,0)
  {Emission\\\footnotesize\itshape results downstream};

% ---- row B: memory suppliers + movers ----
\node[existing, minimum width=4.6cm] (pool) at (4.2,-4.0)
  {Global buffer pool\\\footnotesize\itshape fixed size, allocated at start;\\\footnotesize\itshape buffers recycle, never grow};
\node[existing, minimum width=3.0cm] (heap) at (0.0,-4.0)
  {Process heap\\\footnotesize\itshape grows on demand};
\node[newcomp] (spiller) at (9.0,-4.0)
  {Spiller thread\\\footnotesize\itshape wakes every 100\,ms\\\footnotesize\itshape soft / hard thresholds};
\node[storage, minimum width=4.2cm] (rocksmem) at (14.4,-4.0)
  {RocksDB (in-process)\\\footnotesize\itshape memtable 16\,MB\\\footnotesize\itshape block cache 8\,MB};

% ---- below the cap: the storage device ----
\node[storage, double, double distance=1pt, minimum width=4.2cm] (sd) at (14.4,-8.3)
  {SD card\\\footnotesize\itshape RocksDB spill files};

% ================= GROUP BOX: the memory cap =================
\begin{pgfonlayer}{background}
  \node[draw, dashed, thick, rounded corners, fit=(heap)(pool)(hot)(cold)(spiller)(rocksmem),
        inner xsep=7mm, inner ysep=12mm, fill=black!3] (capbox) {};
\end{pgfonlayer}
\node[anchor=south west, font=\sffamily\footnotesize\bfseries, text=black!70]
  at (capbox.north west) {One engine process under the hard memory cap (cgroup)};
\node[anchor=north west, font=\sffamily\footnotesize\itshape, text=black!60]
  at ([xshift=3mm]capbox.south west) {anonymous memory above the cap $\Rightarrow$ the kernel kills the process (OOM)};

% ---- slice store group box ----
\begin{pgfonlayer}{background}
  \node[draw, thick, rounded corners, fit=(hot)(cold), inner sep=4.5mm, fill=blue!6] (store) {};
\end{pgfonlayer}
\node[anchor=south, font=\sffamily\footnotesize\bfseries, text=blue!40!black]
  at (store.north) {Spillable slice store \textnormal{\footnotesize\itshape (state = hash maps on pages)}};

% ================= EDGES =================
% processing path
\draw[flow] (src) -- node[lbl]{tuples} (workers);
\draw[flow] (workers) -- node[lbl]{writes\\(build)} (hot);
\draw[flow] (hot) -- node[lbl]{watermark\\passes} (cold);
\draw[flow] (cold.east) -- node[lbl, pos=0.5]{restore spilled\\slices, then emit} (emit.west);

% memory supply
\draw[flow] (pool.north) -- node[lbl, pos=0.72]{buffers for\\input and output} (workers.south);
\draw[flow] (pool.north east) -- node[lbl, pos=0.55]{pages\\(bounded allocation)} (hot.south west);
\draw[flow] (heap.north) -- ++(0,0.7) -| node[lbl, pos=0.08]{bucket directories; all pages\\under the 4096-byte default} (hot.south);

% spilling and restoring
% label positions: "watches" sits low near the spiller, "evicts" high near the
% cold slices, so the two boxes occupy separate lanes and cannot collide.
\draw[watch] (spiller.north) -- node[lbl, pos=0.24]{watches live\\slice bytes} (store.south);
\draw[flow] ([xshift=-7mm]cold.south) -- node[lbl, pos=0.30, left]{spiller evicts:\\pages as key--value pairs} ([xshift=-7mm]rocksmem.north);
\draw[flow] ([xshift=-9mm]rocksmem.south) -- node[lbl, pos=0.62]{flush,\\compaction} ([xshift=-9mm]sd.north);
\draw[flow] ([xshift=9mm]sd.north) -- node[lbl, pos=0.38]{read on\\restore} ([xshift=9mm]rocksmem.south);
\draw[flow] ([xshift=7mm]rocksmem.north) -- node[lbl, pos=0.45, right]{restored slices\\(byte copy into pages)} ([xshift=7mm]cold.south);

% ================= LEGEND =================
\node[draw, thick, rounded corners, inner sep=5pt, font=\sffamily\footnotesize]
  at (2.6,-8.3)
  {\tikz\node[draw, thick, fill=black!12, minimum width=9mm, minimum height=4mm, inner sep=0pt]{};\,existing NebulaStream\quad
   \tikz\node[draw, thick, fill=blue!18, minimum width=9mm, minimum height=4mm, inner sep=0pt]{};\,this thesis\quad
   \tikz\node[draw, thick, fill=orange!28, minimum width=9mm, minimum height=4mm, inner sep=0pt]{};\,storage engine / device\quad
   \begin{tikzpicture}\draw[watch](0,0)--(0.8,0);\end{tikzpicture}\,observes\quad
   \begin{tikzpicture}\draw[flow](0,0)--(0.8,0);\end{tikzpicture}\,data / control flow};

\end{tikzpicture}%
}
\caption[System overview: the path of window state through the components]{%
System overview: the path window state takes through the components. Worker
threads write arriving tuples into the hot slice of the current window; its
state lives in hash maps whose pages come from the global buffer pool (bounded
allocation) or, under the engine's 4096-byte default, from the process heap.
Once the watermark passes, the completed window becomes a cold slice and is
held for the emit lag $L$. The spiller thread watches the live slice bytes
against the soft and hard thresholds and evicts cold slices into RocksDB,
which flushes them to the SD card. Before emission, spilled slices are read
back and restored by a plain byte copy into pages, and the results leave the
engine. Everything above the SD card runs in one process under the hard memory
cap; if its anonymous memory exceeds the cap, the kernel kills the process.
Figure~\ref{fig:ooc-integration} shows the same components at the class level.}
\label{fig:system-overview}
\end{figure}

Two thresholds drive it: a soft threshold triggers preemptive spilling, and a hard threshold forces synchronous spilling under acute pressure (Section~\ref{sec:architecture}).
Three concerns are then layered onto this core.
The spiller reacts to the slice budget alone, so a second iteration adds levers that react to pressure on the shared buffer pool (Section~\ref{sec:levers}).
An allocation can still fail when the pool empties, so a third iteration adds allocation backpressure (Section~\ref{sec:backpressure}).
The controller sees a slice only when it is created and never returns spilled memory to the operating system, so a fourth iteration adds a growth signal and page release (Section~\ref{sec:growth-release}).
Every mechanism beyond the core is off by default, so the shipped baseline is byte-identical to the unmodified engine.

\section{Requirements\label{sec:requirements}}

% [5.1]
Four requirements follow from the goal in Chapter~\ref{cha:problem}.
First, the backend spills state before memory pressure becomes fatal: a soft threshold triggers preemptive background spilling so that cold state leaves memory while the system is still healthy.
Second, preemptive spilling can fall behind.
When pressure crosses a hard threshold despite it, the backend forces a synchronous spill on the allocating thread.
The forced spill evicts state as quickly as the storage device permits, so an \ac{OOM} failure is averted rather than merely deferred.
Third, it never changes query results.
A window that was spilled and restored must produce the same output, byte for byte, as one that never left memory; Section~\ref{sec:reproducibility} describes how this is verified.
Fourth, it leaves the just-in-time compilation pipeline untouched.

\section{The Spillable Slice Store\label{sec:architecture}}

% [5.2]
NebulaStream keeps window state in a slice store (Section~\ref{sec:bg-windows}).
The design replaces the in-memory slice store with a spillable slice store behind the same interface.
The operators above the store do not change.

% [5.3]
A slice moves through five states.
A hot slice (the implementation's \texttt{WINDOW\_FILLING} state) belongs to a window that still receives data.
A cold slice belongs to a window that is complete but not yet emitted.
A spilled slice lies in RocksDB.
A restored slice is read back into memory for emission.
An emitted slice has delivered its results and is freed.
Figure~\ref{fig:ooc-integration} shows the store, its new components, and their place in the engine.

% Benjamin's existing TikZ diagram (from the May draft), restored into this project.
% NebulaStream Out-of-Core integration architecture.
% Required preamble (add to Template.tex if missing):
%   \usepackage{tikz}
%   \usetikzlibrary{positioning,arrows.meta,shapes.geometric,fit,backgrounds,calc}
%
% Use as: % NebulaStream Out-of-Core integration architecture.
% Required preamble (add to Template.tex if missing):
%   \usepackage{tikz}
%   \usetikzlibrary{positioning,arrows.meta,shapes.geometric,fit,backgrounds,calc}
%
% Use as: % NebulaStream Out-of-Core integration architecture.
% Required preamble (add to Template.tex if missing):
%   \usepackage{tikz}
%   \usetikzlibrary{positioning,arrows.meta,shapes.geometric,fit,backgrounds,calc}
%
% Use as: \input{figures/spillable_slice_store}
%
% LAYOUT -- strict raster, generous spacing, every edge label boxed.
%   Column centres : x = -1, 6, 13      (pitch 7 cm, gap 3 cm between 4-cm boxes)
%   Build / probe  : x = 4, x = 8       (inside column B, 1 cm gap between 3-cm boxes)
%   Row centres    : y = 0, -2.5, -5, -7.5, -10, -12.5, -15
%                                       (pitch 2.5 cm, gap 1.1 cm between 1.4-cm boxes)
%   Edge labels    : every label rendered as a small white badge (style "lbl")
%                    so it never blurs into an arrow, a fit-box fill, or a node.

\begin{figure}[p]
\centering
\resizebox{\textwidth}{!}{%
\begin{tikzpicture}[
  font=\sffamily\small,
  every node/.style={align=center},
  % ---- node styles ----
  box/.style ={rectangle, draw, thick, minimum width=4cm,   minimum height=1.4cm, inner sep=3pt, rounded corners=1pt},
  sm/.style  ={rectangle, draw, thick, minimum width=3cm,   minimum height=1.6cm, inner sep=3pt, rounded corners=1pt},
  wd/.style  ={rectangle, draw, thick, minimum width=5.5cm, minimum height=1.4cm, inner sep=3pt, rounded corners=1pt},
  existing/.style ={box, fill=black!12},
  newcomp/.style  ={box, fill=blue!18},
  storage/.style  ={box, fill=orange!28},
  smnew/.style    ={sm,  fill=blue!18},
  hook/.style     ={wd,  fill=yellow!55},
  iface/.style    ={wd,  fill=black!12},
  % ---- edge + label styles ----
  flow/.style={-{Latex[length=2.4mm]}, thick},
  sibling/.style={{Latex[length=2.4mm]}-{Latex[length=2.4mm]}, dashed, thick, black!55},
  lbl/.style={
    fill=white, draw=black!30, line width=0.2pt, rounded corners=1pt,
    inner sep=2pt, font=\sffamily\footnotesize\itshape, text=black!75,
  },
]

% ============ grid anchors ============
\def\cA{-1}      \def\cB{6}      \def\cC{13}
\def\cBL{4}      \def\cBR{8}                       % build / probe sub-columns
\def\cMid{2.5}                                     % midpoint of cA and cB

\def\rH{0}       \def\rL{-2.5}                     % header rows
\def\rOne{-5}    \def\rTwo{-7.5}  \def\rThr{-10}   % three state rows
\def\rFou{-12.5} \def\rFiv{-15}                    % backend + disk rows
\def\rLeg{-18}                                     % legend

% ============ HEADER BAND ============
\node[iface]    (iface)   at (\cMid,\rH) {WindowSlicesStoreInterface\\\footnotesize\itshape 8 virtual methods};
\node[hook]     (lower)   at (\cMid,\rL) {LowerToPhysicalWindowedAggregation\\\footnotesize\itshape nes-query-optimizer rewrite rule};
\node[existing] (handler) at (\cC,\rH)   {WindowBasedOperator\-Handler\\::onWatermark()};

% ============ COLUMN A : existing in-memory path ============
\node[existing] (defstore) at (\cA,\rOne) {DefaultTime\-Based\-SliceStore};
\node[existing] (aggslice) at (\cA,\rTwo) {AggregationSlice\\\footnotesize\itshape : HashMapSlice};
\node[existing] (hmaps)    at (\cA,\rThr) {HashMap[0..N$-$1]\\\footnotesize\itshape one per worker thread};
\node[existing] (chained)  at (\cA,\rFou) {ChainedHashMap\\\footnotesize\itshape TupleBuffer-backed pages};

% ============ COLUMN B : new spillable path ============
\node[newcomp] (spillstore)   at (\cB,\rOne)  {SpillableTime\-Based\\SliceStore};
\node[newcomp] (spillslice)   at (\cB,\rTwo)  {SpillableAggregationSlice\\\footnotesize\itshape : HashMapSlice (subclass)};
\node[smnew]   (buildpath)    at (\cBL,\rThr) {Build path\\\footnotesize\itshape getHashMapPtr\\\footnotesize\itshape state == FILLING};
\node[smnew]   (probepath)    at (\cBR,\rThr) {Probe path\\\footnotesize\itshape force-unspill\\\footnotesize\itshape before JIT probe};
\node[newcomp] (backendiface) at (\cB,\rFou)  {SpillBackend\\\footnotesize\itshape interface (RocksDB / file / mmap)};
% ---- THIS THESIS'S EXTENSIONS, called out explicitly ----
% Background spiller thread: sits in the empty midpoint space between the existing
% path and the spill path; drives preemptive/forced spilling off the thresholds.
\node[smnew] (spiller) at (\cMid,\rFiv) {Background\\spiller\\\footnotesize\itshape 100\,ms poll\\\footnotesize\itshape soft / hard\\\footnotesize\itshape thresholds};
% rocks + local disk merged into one storage node to save a row -- a thin
% double rule conveys "physical storage" without needing a second box.
\node[storage, double, double distance=1pt] (rocks) at (\cB,\rFiv)
  {RocksDBSpillBackend\\\footnotesize\itshape local disk};

% ============ COLUMN C : config / memory / watermark ============
\node[newcomp]  (cfg)    at (\cC,\rOne) {SpillConfig\\\footnotesize\itshape extends BaseConfiguration\\\footnotesize\itshape (YAML, nes-configurations)};
\node[newcomp]  (budget) at (\cC,\rTwo) {MemoryBudgetTracker\\\footnotesize\itshape observes BufferManager\\\footnotesize\itshape (advisory, not parallel pool)};
\node[existing] (bufmgr) at (\cC,\rThr) {BufferManager\\\footnotesize\itshape nes-memory\\\footnotesize\itshape\color{blue!60!black}+ availability query\\\footnotesize\itshape\color{blue!60!black}+ pool-pressure hooks};
\node[existing] (wm)     at (\cC,\rFou) {Stream Watermark};

% SpillConfig field detail -- sits next to the rocks row in otherwise empty
% column-C space, two row pitches below wm so it cannot graze anything.
\node[draw, dashed, rounded corners, minimum width=4cm, inner sep=4pt,
      align=left, font=\sffamily\footnotesize] (cfgdetail) at (\cC,\rFiv)
  {\textbf{SpillConfig fields}\\
   \,-- soft / hard memory thresholds\\
   \,-- compression type, level\\
   \,-- level0 slowdown trigger\\
   \,-- write buffer size\\
   \,-- RocksDB path};

% ============ EDGES : header band ============
\draw[flow] (iface) -- node[lbl]{implements} (lower);
\draw[flow] (lower.west)  -- (\cA,\rL)  -- node[lbl, pos=0.55]{creates (default)}  (defstore.north);
\draw[flow] (lower.east)  -- (\cB,\rL)  -- node[lbl, pos=0.55]{creates (spill on)} (spillstore.north);
\draw[flow] (handler) -- node[lbl, pos=0.55]{triggers,\\drives spill} (lower);

% ============ EDGES : existing path (pure vertical, labels offset right) ============
\draw[flow] (defstore) -- node[lbl]{contains} (aggslice);
\draw[flow] (aggslice) -- node[lbl]{owns N}   (hmaps);
\draw[flow] (hmaps)    -- node[lbl]{uses}     (chained);

% ============ EDGES : spillable path ============
\draw[flow] (spillstore) -- node[lbl]{contains} (spillslice);
\draw[flow] (spillslice.south) -- node[lbl, pos=0.55]{write}   (buildpath.north);
\draw[flow] (spillslice.south) -- node[lbl, pos=0.55]{read}    (probepath.north);
\draw[flow] (buildpath.south)  -- node[lbl, pos=0.55]{spill}   (backendiface.north);
\draw[flow] (probepath.south)  -- node[lbl, pos=0.55]{unspill} (backendiface.north);
\draw[flow] (backendiface) -- (rocks);

% ---- spiller thread edges: watches the store budget, drives eviction ----
\draw[flow] (spiller.north) -- node[lbl, pos=0.55]{watches budget} (spillstore.west);
\draw[flow] (spiller.east)  -- node[lbl, pos=0.55]{drives spill} (backendiface.west);

% Sibling relation -- both slice types share the HashMapSlice base.
\draw[sibling] (aggslice.east) -- node[lbl]{shared base} (spillslice.west);

% ============ EDGES : config / memory / watermark ============
\draw[flow] (cfg.west)     -- node[lbl]{tunes}     (spillstore.east);
\draw[flow] (budget.west)  -- node[lbl]{advises}   (spillslice.east);
\draw[flow] (bufmgr.north) -- node[lbl]{observes}  (budget.south);
% Watermark to probe path: diagonal stays left of BufferManager (passes ~x=10.5)
\draw[flow] (wm.west) -- node[lbl, pos=0.4]{end$<$wm $\wedge$ $\neq$FILLING} (probepath.east);

% ============ DASHED GROUP BOXES ============
\begin{pgfonlayer}{background}
  \node[draw, dashed, rounded corners, fit=(defstore)(chained),
        inner sep=6mm, fill=black!3] (existingbox) {};
  \node[draw, dashed, rounded corners, fit=(spillstore)(rocks)(buildpath)(probepath)(spiller),
        inner sep=6mm, fill=blue!4]  (newbox)      {};
\end{pgfonlayer}

% ============ LEGEND : two rows, well below the data plane ============
\node[draw, thick, rounded corners, inner sep=5pt, align=center, font=\sffamily\footnotesize]
  (legend) at ($(\cA,\rLeg)!0.5!(\cC,\rLeg)$)
  {\textbf{Legend}\\[4pt]
   \tikz\node[draw, thick, fill=black!12, minimum width=10mm, minimum height=4mm, inner sep=0pt]{};\,Existing NES code\quad
   \tikz\node[draw, thick, fill=blue!18,  minimum width=10mm, minimum height=4mm, inner sep=0pt]{};\,New component\quad
   \tikz\node[draw, thick, double, double distance=1pt, fill=orange!28, minimum width=10mm, minimum height=4mm, inner sep=0pt]{};\,Persistent storage\quad
   \tikz\node[draw, thick, fill=yellow!55,minimum width=10mm, minimum height=4mm, inner sep=0pt]{};\,Integration point\\[4pt]
   \begin{tikzpicture}\draw[-{Latex[length=2.4mm]}, thick](0,0)--(0.9,0);\end{tikzpicture}\,owns / contains \qquad
   \begin{tikzpicture}\draw[{Latex[length=2.4mm]}-{Latex[length=2.4mm]}, dashed, thick, black!55](0,0)--(0.9,0);\end{tikzpicture}\,shared base class};

\end{tikzpicture}%
}
\caption[NebulaStream out-of-core integration architecture]{%
Integration of the spillable slice store into NebulaStream. Blue nodes are the
components this thesis adds; grey nodes are existing NebulaStream code. The
extensions are the spillable slice store and its slice subclass, the background
spiller thread that drives preemptive and forced spilling off the soft and hard
thresholds, the \texttt{SpillBackend} interface with its RocksDB implementation,
the \texttt{SpillConfig} and \texttt{MemoryBudgetTracker}, and two additions to
the existing \texttt{BufferManager}: an availability query and the pool-pressure
hooks the levers of Section~\ref{sec:levers} use. The optimizer rewrite rule
\texttt{LowerToPhysicalWindowedAggregation} picks the store implementation per
aggregation operator. \texttt{SpillableAggregationSlice} is a \emph{subclass} of
\texttt{HashMapSlice} (not a wrapper), so the Nautilus just-in-time (JIT) build path can keep
using \texttt{getHashMapPtr(WorkerId)} unchanged; only the probe path is aware of
unspill. The probe path must force-unspill \emph{before} the compiled probe
pipeline runs: the query compiler emits machine code that reads the hash map
through raw pointers, so the state must be resident in memory before that code
executes, and the compiled code itself stays unchanged. A slice holds one hashmap per worker thread and becomes eligible for
spilling once its end timestamp falls below the watermark and its state has left
\texttt{WINDOW\_FILLING}. The \texttt{SpillBackend} interface keeps the RocksDB
backend swappable for evaluation. \texttt{MemoryBudgetTracker} observes the
existing \texttt{BufferManager} rather than maintaining a parallel memory pool.}
\label{fig:ooc-integration}
\end{figure}

%
% LAYOUT -- strict raster, generous spacing, every edge label boxed.
%   Column centres : x = -1, 6, 13      (pitch 7 cm, gap 3 cm between 4-cm boxes)
%   Build / probe  : x = 4, x = 8       (inside column B, 1 cm gap between 3-cm boxes)
%   Row centres    : y = 0, -2.5, -5, -7.5, -10, -12.5, -15
%                                       (pitch 2.5 cm, gap 1.1 cm between 1.4-cm boxes)
%   Edge labels    : every label rendered as a small white badge (style "lbl")
%                    so it never blurs into an arrow, a fit-box fill, or a node.

\begin{figure}[p]
\centering
\resizebox{\textwidth}{!}{%
\begin{tikzpicture}[
  font=\sffamily\small,
  every node/.style={align=center},
  % ---- node styles ----
  box/.style ={rectangle, draw, thick, minimum width=4cm,   minimum height=1.4cm, inner sep=3pt, rounded corners=1pt},
  sm/.style  ={rectangle, draw, thick, minimum width=3cm,   minimum height=1.6cm, inner sep=3pt, rounded corners=1pt},
  wd/.style  ={rectangle, draw, thick, minimum width=5.5cm, minimum height=1.4cm, inner sep=3pt, rounded corners=1pt},
  existing/.style ={box, fill=black!12},
  newcomp/.style  ={box, fill=blue!18},
  storage/.style  ={box, fill=orange!28},
  smnew/.style    ={sm,  fill=blue!18},
  hook/.style     ={wd,  fill=yellow!55},
  iface/.style    ={wd,  fill=black!12},
  % ---- edge + label styles ----
  flow/.style={-{Latex[length=2.4mm]}, thick},
  sibling/.style={{Latex[length=2.4mm]}-{Latex[length=2.4mm]}, dashed, thick, black!55},
  lbl/.style={
    fill=white, draw=black!30, line width=0.2pt, rounded corners=1pt,
    inner sep=2pt, font=\sffamily\footnotesize\itshape, text=black!75,
  },
]

% ============ grid anchors ============
\def\cA{-1}      \def\cB{6}      \def\cC{13}
\def\cBL{4}      \def\cBR{8}                       % build / probe sub-columns
\def\cMid{2.5}                                     % midpoint of cA and cB

\def\rH{0}       \def\rL{-2.5}                     % header rows
\def\rOne{-5}    \def\rTwo{-7.5}  \def\rThr{-10}   % three state rows
\def\rFou{-12.5} \def\rFiv{-15}                    % backend + disk rows
\def\rLeg{-18}                                     % legend

% ============ HEADER BAND ============
\node[iface]    (iface)   at (\cMid,\rH) {WindowSlicesStoreInterface\\\footnotesize\itshape 8 virtual methods};
\node[hook]     (lower)   at (\cMid,\rL) {LowerToPhysicalWindowedAggregation\\\footnotesize\itshape nes-query-optimizer rewrite rule};
\node[existing] (handler) at (\cC,\rH)   {WindowBasedOperator\-Handler\\::onWatermark()};

% ============ COLUMN A : existing in-memory path ============
\node[existing] (defstore) at (\cA,\rOne) {DefaultTime\-Based\-SliceStore};
\node[existing] (aggslice) at (\cA,\rTwo) {AggregationSlice\\\footnotesize\itshape : HashMapSlice};
\node[existing] (hmaps)    at (\cA,\rThr) {HashMap[0..N$-$1]\\\footnotesize\itshape one per worker thread};
\node[existing] (chained)  at (\cA,\rFou) {ChainedHashMap\\\footnotesize\itshape TupleBuffer-backed pages};

% ============ COLUMN B : new spillable path ============
\node[newcomp] (spillstore)   at (\cB,\rOne)  {SpillableTime\-Based\\SliceStore};
\node[newcomp] (spillslice)   at (\cB,\rTwo)  {SpillableAggregationSlice\\\footnotesize\itshape : HashMapSlice (subclass)};
\node[smnew]   (buildpath)    at (\cBL,\rThr) {Build path\\\footnotesize\itshape getHashMapPtr\\\footnotesize\itshape state == FILLING};
\node[smnew]   (probepath)    at (\cBR,\rThr) {Probe path\\\footnotesize\itshape force-unspill\\\footnotesize\itshape before JIT probe};
\node[newcomp] (backendiface) at (\cB,\rFou)  {SpillBackend\\\footnotesize\itshape interface (RocksDB / file / mmap)};
% ---- THIS THESIS'S EXTENSIONS, called out explicitly ----
% Background spiller thread: sits in the empty midpoint space between the existing
% path and the spill path; drives preemptive/forced spilling off the thresholds.
\node[smnew] (spiller) at (\cMid,\rFiv) {Background\\spiller\\\footnotesize\itshape 100\,ms poll\\\footnotesize\itshape soft / hard\\\footnotesize\itshape thresholds};
% rocks + local disk merged into one storage node to save a row -- a thin
% double rule conveys "physical storage" without needing a second box.
\node[storage, double, double distance=1pt] (rocks) at (\cB,\rFiv)
  {RocksDBSpillBackend\\\footnotesize\itshape local disk};

% ============ COLUMN C : config / memory / watermark ============
\node[newcomp]  (cfg)    at (\cC,\rOne) {SpillConfig\\\footnotesize\itshape extends BaseConfiguration\\\footnotesize\itshape (YAML, nes-configurations)};
\node[newcomp]  (budget) at (\cC,\rTwo) {MemoryBudgetTracker\\\footnotesize\itshape observes BufferManager\\\footnotesize\itshape (advisory, not parallel pool)};
\node[existing] (bufmgr) at (\cC,\rThr) {BufferManager\\\footnotesize\itshape nes-memory\\\footnotesize\itshape\color{blue!60!black}+ availability query\\\footnotesize\itshape\color{blue!60!black}+ pool-pressure hooks};
\node[existing] (wm)     at (\cC,\rFou) {Stream Watermark};

% SpillConfig field detail -- sits next to the rocks row in otherwise empty
% column-C space, two row pitches below wm so it cannot graze anything.
\node[draw, dashed, rounded corners, minimum width=4cm, inner sep=4pt,
      align=left, font=\sffamily\footnotesize] (cfgdetail) at (\cC,\rFiv)
  {\textbf{SpillConfig fields}\\
   \,-- soft / hard memory thresholds\\
   \,-- compression type, level\\
   \,-- level0 slowdown trigger\\
   \,-- write buffer size\\
   \,-- RocksDB path};

% ============ EDGES : header band ============
\draw[flow] (iface) -- node[lbl]{implements} (lower);
\draw[flow] (lower.west)  -- (\cA,\rL)  -- node[lbl, pos=0.55]{creates (default)}  (defstore.north);
\draw[flow] (lower.east)  -- (\cB,\rL)  -- node[lbl, pos=0.55]{creates (spill on)} (spillstore.north);
\draw[flow] (handler) -- node[lbl, pos=0.55]{triggers,\\drives spill} (lower);

% ============ EDGES : existing path (pure vertical, labels offset right) ============
\draw[flow] (defstore) -- node[lbl]{contains} (aggslice);
\draw[flow] (aggslice) -- node[lbl]{owns N}   (hmaps);
\draw[flow] (hmaps)    -- node[lbl]{uses}     (chained);

% ============ EDGES : spillable path ============
\draw[flow] (spillstore) -- node[lbl]{contains} (spillslice);
\draw[flow] (spillslice.south) -- node[lbl, pos=0.55]{write}   (buildpath.north);
\draw[flow] (spillslice.south) -- node[lbl, pos=0.55]{read}    (probepath.north);
\draw[flow] (buildpath.south)  -- node[lbl, pos=0.55]{spill}   (backendiface.north);
\draw[flow] (probepath.south)  -- node[lbl, pos=0.55]{unspill} (backendiface.north);
\draw[flow] (backendiface) -- (rocks);

% ---- spiller thread edges: watches the store budget, drives eviction ----
\draw[flow] (spiller.north) -- node[lbl, pos=0.55]{watches budget} (spillstore.west);
\draw[flow] (spiller.east)  -- node[lbl, pos=0.55]{drives spill} (backendiface.west);

% Sibling relation -- both slice types share the HashMapSlice base.
\draw[sibling] (aggslice.east) -- node[lbl]{shared base} (spillslice.west);

% ============ EDGES : config / memory / watermark ============
\draw[flow] (cfg.west)     -- node[lbl]{tunes}     (spillstore.east);
\draw[flow] (budget.west)  -- node[lbl]{advises}   (spillslice.east);
\draw[flow] (bufmgr.north) -- node[lbl]{observes}  (budget.south);
% Watermark to probe path: diagonal stays left of BufferManager (passes ~x=10.5)
\draw[flow] (wm.west) -- node[lbl, pos=0.4]{end$<$wm $\wedge$ $\neq$FILLING} (probepath.east);

% ============ DASHED GROUP BOXES ============
\begin{pgfonlayer}{background}
  \node[draw, dashed, rounded corners, fit=(defstore)(chained),
        inner sep=6mm, fill=black!3] (existingbox) {};
  \node[draw, dashed, rounded corners, fit=(spillstore)(rocks)(buildpath)(probepath)(spiller),
        inner sep=6mm, fill=blue!4]  (newbox)      {};
\end{pgfonlayer}

% ============ LEGEND : two rows, well below the data plane ============
\node[draw, thick, rounded corners, inner sep=5pt, align=center, font=\sffamily\footnotesize]
  (legend) at ($(\cA,\rLeg)!0.5!(\cC,\rLeg)$)
  {\textbf{Legend}\\[4pt]
   \tikz\node[draw, thick, fill=black!12, minimum width=10mm, minimum height=4mm, inner sep=0pt]{};\,Existing NES code\quad
   \tikz\node[draw, thick, fill=blue!18,  minimum width=10mm, minimum height=4mm, inner sep=0pt]{};\,New component\quad
   \tikz\node[draw, thick, double, double distance=1pt, fill=orange!28, minimum width=10mm, minimum height=4mm, inner sep=0pt]{};\,Persistent storage\quad
   \tikz\node[draw, thick, fill=yellow!55,minimum width=10mm, minimum height=4mm, inner sep=0pt]{};\,Integration point\\[4pt]
   \begin{tikzpicture}\draw[-{Latex[length=2.4mm]}, thick](0,0)--(0.9,0);\end{tikzpicture}\,owns / contains \qquad
   \begin{tikzpicture}\draw[{Latex[length=2.4mm]}-{Latex[length=2.4mm]}, dashed, thick, black!55](0,0)--(0.9,0);\end{tikzpicture}\,shared base class};

\end{tikzpicture}%
}
\caption[NebulaStream out-of-core integration architecture]{%
Integration of the spillable slice store into NebulaStream. Blue nodes are the
components this thesis adds; grey nodes are existing NebulaStream code. The
extensions are the spillable slice store and its slice subclass, the background
spiller thread that drives preemptive and forced spilling off the soft and hard
thresholds, the \texttt{SpillBackend} interface with its RocksDB implementation,
the \texttt{SpillConfig} and \texttt{MemoryBudgetTracker}, and two additions to
the existing \texttt{BufferManager}: an availability query and the pool-pressure
hooks the levers of Section~\ref{sec:levers} use. The optimizer rewrite rule
\texttt{LowerToPhysicalWindowedAggregation} picks the store implementation per
aggregation operator. \texttt{SpillableAggregationSlice} is a \emph{subclass} of
\texttt{HashMapSlice} (not a wrapper), so the Nautilus just-in-time (JIT) build path can keep
using \texttt{getHashMapPtr(WorkerId)} unchanged; only the probe path is aware of
unspill. The probe path must force-unspill \emph{before} the compiled probe
pipeline runs: the query compiler emits machine code that reads the hash map
through raw pointers, so the state must be resident in memory before that code
executes, and the compiled code itself stays unchanged. A slice holds one hashmap per worker thread and becomes eligible for
spilling once its end timestamp falls below the watermark and its state has left
\texttt{WINDOW\_FILLING}. The \texttt{SpillBackend} interface keeps the RocksDB
backend swappable for evaluation. \texttt{MemoryBudgetTracker} observes the
existing \texttt{BufferManager} rather than maintaining a parallel memory pool.}
\label{fig:ooc-integration}
\end{figure}

%
% LAYOUT -- strict raster, generous spacing, every edge label boxed.
%   Column centres : x = -1, 6, 13      (pitch 7 cm, gap 3 cm between 4-cm boxes)
%   Build / probe  : x = 4, x = 8       (inside column B, 1 cm gap between 3-cm boxes)
%   Row centres    : y = 0, -2.5, -5, -7.5, -10, -12.5, -15
%                                       (pitch 2.5 cm, gap 1.1 cm between 1.4-cm boxes)
%   Edge labels    : every label rendered as a small white badge (style "lbl")
%                    so it never blurs into an arrow, a fit-box fill, or a node.

\begin{figure}[p]
\centering
\resizebox{\textwidth}{!}{%
\begin{tikzpicture}[
  font=\sffamily\small,
  every node/.style={align=center},
  % ---- node styles ----
  box/.style ={rectangle, draw, thick, minimum width=4cm,   minimum height=1.4cm, inner sep=3pt, rounded corners=1pt},
  sm/.style  ={rectangle, draw, thick, minimum width=3cm,   minimum height=1.6cm, inner sep=3pt, rounded corners=1pt},
  wd/.style  ={rectangle, draw, thick, minimum width=5.5cm, minimum height=1.4cm, inner sep=3pt, rounded corners=1pt},
  existing/.style ={box, fill=black!12},
  newcomp/.style  ={box, fill=blue!18},
  storage/.style  ={box, fill=orange!28},
  smnew/.style    ={sm,  fill=blue!18},
  hook/.style     ={wd,  fill=yellow!55},
  iface/.style    ={wd,  fill=black!12},
  % ---- edge + label styles ----
  flow/.style={-{Latex[length=2.4mm]}, thick},
  sibling/.style={{Latex[length=2.4mm]}-{Latex[length=2.4mm]}, dashed, thick, black!55},
  lbl/.style={
    fill=white, draw=black!30, line width=0.2pt, rounded corners=1pt,
    inner sep=2pt, font=\sffamily\footnotesize\itshape, text=black!75,
  },
]

% ============ grid anchors ============
\def\cA{-1}      \def\cB{6}      \def\cC{13}
\def\cBL{4}      \def\cBR{8}                       % build / probe sub-columns
\def\cMid{2.5}                                     % midpoint of cA and cB

\def\rH{0}       \def\rL{-2.5}                     % header rows
\def\rOne{-5}    \def\rTwo{-7.5}  \def\rThr{-10}   % three state rows
\def\rFou{-12.5} \def\rFiv{-15}                    % backend + disk rows
\def\rLeg{-18}                                     % legend

% ============ HEADER BAND ============
\node[iface]    (iface)   at (\cMid,\rH) {WindowSlicesStoreInterface\\\footnotesize\itshape 8 virtual methods};
\node[hook]     (lower)   at (\cMid,\rL) {LowerToPhysicalWindowedAggregation\\\footnotesize\itshape nes-query-optimizer rewrite rule};
\node[existing] (handler) at (\cC,\rH)   {WindowBasedOperator\-Handler\\::onWatermark()};

% ============ COLUMN A : existing in-memory path ============
\node[existing] (defstore) at (\cA,\rOne) {DefaultTime\-Based\-SliceStore};
\node[existing] (aggslice) at (\cA,\rTwo) {AggregationSlice\\\footnotesize\itshape : HashMapSlice};
\node[existing] (hmaps)    at (\cA,\rThr) {HashMap[0..N$-$1]\\\footnotesize\itshape one per worker thread};
\node[existing] (chained)  at (\cA,\rFou) {ChainedHashMap\\\footnotesize\itshape TupleBuffer-backed pages};

% ============ COLUMN B : new spillable path ============
\node[newcomp] (spillstore)   at (\cB,\rOne)  {SpillableTime\-Based\\SliceStore};
\node[newcomp] (spillslice)   at (\cB,\rTwo)  {SpillableAggregationSlice\\\footnotesize\itshape : HashMapSlice (subclass)};
\node[smnew]   (buildpath)    at (\cBL,\rThr) {Build path\\\footnotesize\itshape getHashMapPtr\\\footnotesize\itshape state == FILLING};
\node[smnew]   (probepath)    at (\cBR,\rThr) {Probe path\\\footnotesize\itshape force-unspill\\\footnotesize\itshape before JIT probe};
\node[newcomp] (backendiface) at (\cB,\rFou)  {SpillBackend\\\footnotesize\itshape interface (RocksDB / file / mmap)};
% ---- THIS THESIS'S EXTENSIONS, called out explicitly ----
% Background spiller thread: sits in the empty midpoint space between the existing
% path and the spill path; drives preemptive/forced spilling off the thresholds.
\node[smnew] (spiller) at (\cMid,\rFiv) {Background\\spiller\\\footnotesize\itshape 100\,ms poll\\\footnotesize\itshape soft / hard\\\footnotesize\itshape thresholds};
% rocks + local disk merged into one storage node to save a row -- a thin
% double rule conveys "physical storage" without needing a second box.
\node[storage, double, double distance=1pt] (rocks) at (\cB,\rFiv)
  {RocksDBSpillBackend\\\footnotesize\itshape local disk};

% ============ COLUMN C : config / memory / watermark ============
\node[newcomp]  (cfg)    at (\cC,\rOne) {SpillConfig\\\footnotesize\itshape extends BaseConfiguration\\\footnotesize\itshape (YAML, nes-configurations)};
\node[newcomp]  (budget) at (\cC,\rTwo) {MemoryBudgetTracker\\\footnotesize\itshape observes BufferManager\\\footnotesize\itshape (advisory, not parallel pool)};
\node[existing] (bufmgr) at (\cC,\rThr) {BufferManager\\\footnotesize\itshape nes-memory\\\footnotesize\itshape\color{blue!60!black}+ availability query\\\footnotesize\itshape\color{blue!60!black}+ pool-pressure hooks};
\node[existing] (wm)     at (\cC,\rFou) {Stream Watermark};

% SpillConfig field detail -- sits next to the rocks row in otherwise empty
% column-C space, two row pitches below wm so it cannot graze anything.
\node[draw, dashed, rounded corners, minimum width=4cm, inner sep=4pt,
      align=left, font=\sffamily\footnotesize] (cfgdetail) at (\cC,\rFiv)
  {\textbf{SpillConfig fields}\\
   \,-- soft / hard memory thresholds\\
   \,-- compression type, level\\
   \,-- level0 slowdown trigger\\
   \,-- write buffer size\\
   \,-- RocksDB path};

% ============ EDGES : header band ============
\draw[flow] (iface) -- node[lbl]{implements} (lower);
\draw[flow] (lower.west)  -- (\cA,\rL)  -- node[lbl, pos=0.55]{creates (default)}  (defstore.north);
\draw[flow] (lower.east)  -- (\cB,\rL)  -- node[lbl, pos=0.55]{creates (spill on)} (spillstore.north);
\draw[flow] (handler) -- node[lbl, pos=0.55]{triggers,\\drives spill} (lower);

% ============ EDGES : existing path (pure vertical, labels offset right) ============
\draw[flow] (defstore) -- node[lbl]{contains} (aggslice);
\draw[flow] (aggslice) -- node[lbl]{owns N}   (hmaps);
\draw[flow] (hmaps)    -- node[lbl]{uses}     (chained);

% ============ EDGES : spillable path ============
\draw[flow] (spillstore) -- node[lbl]{contains} (spillslice);
\draw[flow] (spillslice.south) -- node[lbl, pos=0.55]{write}   (buildpath.north);
\draw[flow] (spillslice.south) -- node[lbl, pos=0.55]{read}    (probepath.north);
\draw[flow] (buildpath.south)  -- node[lbl, pos=0.55]{spill}   (backendiface.north);
\draw[flow] (probepath.south)  -- node[lbl, pos=0.55]{unspill} (backendiface.north);
\draw[flow] (backendiface) -- (rocks);

% ---- spiller thread edges: watches the store budget, drives eviction ----
\draw[flow] (spiller.north) -- node[lbl, pos=0.55]{watches budget} (spillstore.west);
\draw[flow] (spiller.east)  -- node[lbl, pos=0.55]{drives spill} (backendiface.west);

% Sibling relation -- both slice types share the HashMapSlice base.
\draw[sibling] (aggslice.east) -- node[lbl]{shared base} (spillslice.west);

% ============ EDGES : config / memory / watermark ============
\draw[flow] (cfg.west)     -- node[lbl]{tunes}     (spillstore.east);
\draw[flow] (budget.west)  -- node[lbl]{advises}   (spillslice.east);
\draw[flow] (bufmgr.north) -- node[lbl]{observes}  (budget.south);
% Watermark to probe path: diagonal stays left of BufferManager (passes ~x=10.5)
\draw[flow] (wm.west) -- node[lbl, pos=0.4]{end$<$wm $\wedge$ $\neq$FILLING} (probepath.east);

% ============ DASHED GROUP BOXES ============
\begin{pgfonlayer}{background}
  \node[draw, dashed, rounded corners, fit=(defstore)(chained),
        inner sep=6mm, fill=black!3] (existingbox) {};
  \node[draw, dashed, rounded corners, fit=(spillstore)(rocks)(buildpath)(probepath)(spiller),
        inner sep=6mm, fill=blue!4]  (newbox)      {};
\end{pgfonlayer}

% ============ LEGEND : two rows, well below the data plane ============
\node[draw, thick, rounded corners, inner sep=5pt, align=center, font=\sffamily\footnotesize]
  (legend) at ($(\cA,\rLeg)!0.5!(\cC,\rLeg)$)
  {\textbf{Legend}\\[4pt]
   \tikz\node[draw, thick, fill=black!12, minimum width=10mm, minimum height=4mm, inner sep=0pt]{};\,Existing NES code\quad
   \tikz\node[draw, thick, fill=blue!18,  minimum width=10mm, minimum height=4mm, inner sep=0pt]{};\,New component\quad
   \tikz\node[draw, thick, double, double distance=1pt, fill=orange!28, minimum width=10mm, minimum height=4mm, inner sep=0pt]{};\,Persistent storage\quad
   \tikz\node[draw, thick, fill=yellow!55,minimum width=10mm, minimum height=4mm, inner sep=0pt]{};\,Integration point\\[4pt]
   \begin{tikzpicture}\draw[-{Latex[length=2.4mm]}, thick](0,0)--(0.9,0);\end{tikzpicture}\,owns / contains \qquad
   \begin{tikzpicture}\draw[{Latex[length=2.4mm]}-{Latex[length=2.4mm]}, dashed, thick, black!55](0,0)--(0.9,0);\end{tikzpicture}\,shared base class};

\end{tikzpicture}%
}
\caption[NebulaStream out-of-core integration architecture]{%
Integration of the spillable slice store into NebulaStream. Blue nodes are the
components this thesis adds; grey nodes are existing NebulaStream code. The
extensions are the spillable slice store and its slice subclass, the background
spiller thread that drives preemptive and forced spilling off the soft and hard
thresholds, the \texttt{SpillBackend} interface with its RocksDB implementation,
the \texttt{SpillConfig} and \texttt{MemoryBudgetTracker}, and two additions to
the existing \texttt{BufferManager}: an availability query and the pool-pressure
hooks the levers of Section~\ref{sec:levers} use. The optimizer rewrite rule
\texttt{LowerToPhysicalWindowedAggregation} picks the store implementation per
aggregation operator. \texttt{SpillableAggregationSlice} is a \emph{subclass} of
\texttt{HashMapSlice} (not a wrapper), so the Nautilus just-in-time (JIT) build path can keep
using \texttt{getHashMapPtr(WorkerId)} unchanged; only the probe path is aware of
unspill. The probe path must force-unspill \emph{before} the compiled probe
pipeline runs: the query compiler emits machine code that reads the hash map
through raw pointers, so the state must be resident in memory before that code
executes, and the compiled code itself stays unchanged. A slice holds one hashmap per worker thread and becomes eligible for
spilling once its end timestamp falls below the watermark and its state has left
\texttt{WINDOW\_FILLING}. The \texttt{SpillBackend} interface keeps the RocksDB
backend swappable for evaluation. \texttt{MemoryBudgetTracker} observes the
existing \texttt{BufferManager} rather than maintaining a parallel memory pool.}
\label{fig:ooc-integration}
\end{figure}


% [5.3b — emit lag: the thesis's addition that creates the cold-slice population]
Standard NebulaStream emits a completed window without delay.
This backend instead holds it for a configured emit lag $L$.
During the lag no operator writes to the window anymore, so the spiller has time to move it to storage.
With window length $W$, the store then holds $N_w = L/W$ completed windows as cold slices at any time.

% [5.4]
A background thread, the spiller, wakes every 100 milliseconds.
The soft threshold is the memory level that triggers preemptive spilling.
When the live slice bytes exceed the soft threshold, the spiller moves cold slices to RocksDB.
It moves cold slices first because no operator writes to them anymore.

% [5.5]
The hard threshold is a higher level for forced spilling.
When the live slice bytes reach it, the store spills synchronously.
This covers the case where the background spiller falls behind the arrival rate.

% [5.6]
Slices serialize page by page into RocksDB key-value pairs.
The serializer preserves the exact memory layout of the hash map.
A restore is therefore a plain byte copy back into pages.

% [5.7]-
Before a window's results are emitted, the store restores its spilled slices.

% [5.8]
The spill controller observes one signal: the live byte count of slice data pages.
Other consumers of the buffer pool, such as input buffers and emission buffers, are invisible to it.
This is a design property, stated here openly: the controller defends the slice budget, not the whole pool.

\section{Choosing the RocksDB Backend Configuration\label{sec:rocksdb-config}}

% Ch.4 keeps only the JUSTIFIED SELECTION of the backend knobs. The db_bench
% evaluation (10,000-run sweep + results) now opens Chapter 5 (Section~\ref{sec:rocksdb-eval}).
The store spills slices into RocksDB (Section~\ref{sec:architecture}).
RocksDB exposes many tuning knobs, and prior work shows that three of them govern its behavior under a tight memory budget: the compression algorithm, the write buffer size, and the block cache size.
Compression trades processor time for fewer bytes written to disk~\cite{jin2022}.
The write buffer is the in-memory region RocksDB fills before flushing to disk; a larger buffer flushes less often but commits more memory that then competes with the window state~\cite{rocksdb-memory}.
The block cache holds recently read blocks in memory and speeds up reads at the same memory cost~\cite{rocksdb-memory}.
Based on these findings, this thesis fixes one backend configuration rather than searching the space per experiment.
Table~\ref{tab:rocksdb-config} lists the three knobs, the reason each matters on an edge budget, and the value chosen.

\begin{table}[ht]
\centering
\small
\begin{tabular}{|l|p{8.2cm}|l|}
\hline
\textbf{Knob} & \textbf{Why it matters on an edge budget} & \textbf{Value} \\
\hline
Compression & fewer bytes written to the storage device per spilled slice, at a processor cost~\cite{jin2022} & lz4 \\
\hline
Write buffer & larger flushes less often but commits memory that competes with the window state~\cite{rocksdb-memory} & 16\,MB \\
\hline
Block cache & speeds up reads at the same memory cost; the spill duty cycle is write-heavy~\cite{rocksdb-memory} & 8\,MB \\
\hline
\end{tabular}
\caption{The three RocksDB knobs that govern behavior under a tight memory
budget, and the value chosen for each. The values are small so that the backend
leaves the most memory for window state at the tightest 512\,MB cap. Section~\ref{sec:rocksdb-eval}
reports the sweep that produced these values, and Section~\ref{sec:rb32} tests
whether a larger backend helps once the cap grows.}
\label{tab:rocksdb-config}
\end{table}

The three values are the smallest that still perform well under the tightest 512\,MB cap.
A small write buffer and a small block cache leave the most memory for window state, which is the scarce resource on an edge device.
The evaluation holds this configuration constant across all caps (Section~\ref{sec:setup}).
It also measures whether a larger backend would help once a larger cap leaves room for it (Section~\ref{sec:rb32}).

\section{Design Iteration Two: Pool-Pressure Levers\label{sec:levers}}

% [5.9]
The core spiller watches the slice budget alone.
It stays blind to the shared buffer pool that every other part of the engine draws from.
A run can therefore exhaust the pool and abort while the slice budget still looks healthy, because input buffers, queues, and emission buffers consumed the pool instead.
The second iteration adds four levers that react to pool pressure directly.
Each one is gated behind its own configuration flag.
Lever 1a lets the spiller also drain when the free buffers in the pool fall below a reserve.
Lever 1b reacts to an allocation that finds the pool empty: it triggers an immediate spill of one cold slice instead of waiting passively.
Lever 1c paces restores: a restore before emission proceeds only when the pool has headroom.
Lever 1d adds the bucket directories (Section~\ref{sec:bg-windows}) to the live slice budget.
Whether these levers convert failing runs into completing ones is an empirical question; Section~\ref{sec:rq3} answers it.

\section{Design Iteration Three: Allocation Backpressure\label{sec:backpressure}}

% [5.10]
The spiller can fall behind the arrival rate.
When it does and all buffers are in flight, NebulaStream aborts the next allocation (Section~\ref{sec:bg-nes}).

Stream engines handle this demand--supply gap with backpressure.
Reactive Streams propagate demand signals upstream; Flink uses credit-based flow control; Storm and Heron throttle at the source~\cite{reactivestreams,flink_credit,storm886,kulkarni2015heron}.
This design applies backpressure at the allocation layer: an allocation that would abort stalls instead, holding its thread in bounded rounds while the spiller frees space.

The stall must be bounded.
A progress guard tracks how many buffers return to the pool each round.
If none come back across two consecutive rounds, the threads that would return buffers are blocked on the same pool, so waiting cannot break the cycle.
The guard aborts.

% [5.11]
Every mechanism in this chapter is off by default.
With all gates closed, the engine is byte-identical to the unmodified baseline.
Every comparison in Chapter~\ref{cha:evaluation} is therefore an A/B test on one binary.

\section{Design Iteration Four: The Growth Signal and Page Release\label{sec:growth-release}}

% [5.12]
Two gaps survive the first three iterations.
The controller checks the live budget only when a slice is created.
A high-cardinality aggregation grows one slice between window boundaries without creating a new one, so that growth is unseen until the next safety-net spiller wake.
And a spilled slice frees its logical state but not its resident memory.
The pool is one fixed slab reserved at startup, so a recycled buffer only returns to the free list; the anonymous pages it touched were faulted in once and are never returned to the operating system.

% [5.13]
The growth signal closes the first gap.
It re-evaluates the same soft and hard thresholds on the existing-slice write path, sampled every $N$ writes so the budget read never runs per record.
Placement is the difficulty.
The per-worker slice cache, an engine default (Section~\ref{sec:geometry-valley}), serves steady-state writes without re-entering the store, so a check inside the store's lookup is reached only on a cache miss, the path the growth never takes.
The signal is therefore placed on the data-structure lookup the build operator calls for every record, whether the slice cache hits or misses.
A typical stride $N$ is a few thousand writes; $N = 0$ disables the signal and restores the byte-identical baseline.

% [5.14]
Page release closes the second gap.
On recycle the buffer manager advises the kernel to drop a recycled buffer's physical pages, and does so only while the free list already holds a configured fraction of the pool.
This free-list gate (a pool-occupancy condition, unrelated to the event-time watermark of Section~\ref{sec:bg-windows}) confines the release to the idle, post-drain state and spares the fast-cycling input buffers.
The advice is issued before the buffer re-enters the free list, so no concurrent allocation can fault the page underneath it.
The default sixty-four-byte buffer alignment is too coarse for a per-buffer page operation, so the flag also switches the pool to a page-aligned layout that keeps the control blocks in a separate region.
Both mechanisms are off by default, and page release extends the buffer manager rather than the store.

\section{Integration into NebulaStream\label{sec:integration}}

% [6.0 — provenance: the upstream state all changes build on]
The implementation builds on a public fork of the NebulaStream repository.
The base state is commit \texttt{cc754a94} of the fork's \texttt{main} branch, taken on May 20, 2026, before any change of this thesis.
This base was verified clean: the unmodified tree builds and passes all 752 unit tests.
Every mechanism in this chapter is a change against that state, and the byte-identical baseline of Section~\ref{sec:backpressure} refers to it.

% [6.1]
The store is created when the query is lowered into its physical plan.
Its configuration travels from the worker configuration through the engine builder into the store and the buffer manager.

% [6.2]
One comparison inside ChainedHashMap decides whether a hash map page is pooled.
A page is pooled if the configured page size fits into the engine's buffer size.
The engine's default buffer size is 4096 bytes.
The page size in this thesis is 8192 bytes.
The default therefore keeps every hash map page unpooled, on the process heap rather than in the global buffer pool.

% [6.3]
The buffer manager gains three extensions.
An availability query reports the number of free buffers.
Pressure listeners notify the store when an allocation finds the pool empty; lever 1b uses this hook.
The backpressure loop uses a recycle counter as its progress signal.

\section{Instrumentation\label{sec:instrumentation}}

% [6.4]
Measuring a spilling system cleanly requires separating what the kernel does for bookkeeping from what the workload actually uses.
A spilling backend writes slices to disk and reads them back on restore; the kernel's response is to grow the page cache.
An early evaluation round read undifferentiated \texttt{memory.current} totals and produced footprint numbers that moved with the spill volume rather than with the window working set.
The fix is to read \texttt{memory.stat} instead, which breaks cgroup consumption into anonymous memory, page cache, and shared memory (Section~\ref{sec:bg-metrics}).
A sampler does this every 0.2 seconds; the per-process anonymous set from the kernel's per-process status file runs alongside it as a cross-check on the cgroup accounting.

Five further gauges observe what the application sees.
A pool gauge records occupancy at every buffer allocation, exposing when the pool runs dry and by how much.
An unpooled gauge tracks the live bytes held outside the pool, bounding the share of memory invisible to the spill controller.
A map gauge captures the exact page-count composition of every hash map at its destruction, so the evaluation can attribute footprint to individual window slots.
A spill gauge records the controller's state at every 100-millisecond spiller wake.
A backpressure gauge logs each stall with its duration.

Together, the six probes let Chapter~\ref{cha:evaluation} separate working-set footprint from spill artifacts, trace the controller's decisions under pressure, and cross-examine each claim against independent signals.

\section{Evaluation Harness and Reproducibility\label{sec:reproducibility}}

% [6.5]
The experiments ran on a desktop machine (Intel Core i7-10700K: eight physical cores, sixteen logical threads; 32\,GB RAM) rather than a real edge node.
Docker constrains the container to an edge-class resource envelope; it does not reproduce any specific device.
The memory cap is the central constraint and the only one varied across experiments: a cgroup hard limit with swap disabled, set to 512, 1024, or 2048\,MB, so the backend is held to a fixed memory budget it may never exceed.
Docker imposes this cap through the Linux kernel's control-group (cgroup) memory controller: the \texttt{--memory} flag of Table~\ref{tab:edge-profile} sets the cgroup's hard limit.
The ``Docker memory limit'' and the ``cgroup cap'' named throughout this thesis are therefore the same mechanism seen at the container and kernel layers, not two independent limits, and the memory probes read their footprint from that same cgroup (Section~\ref{sec:bg-metrics}).
The 512\,MB floor is chosen to sit at the low end of the memory budgets reported for constrained \ac{IoT} and edge-class devices~\cite{ojo2019,shi2016}, the regime where in-memory state no longer fits and out-of-core processing becomes necessary; the larger caps trace how the behavior changes as the budget grows.
Two further levers complete the envelope: a two-core CPU budget pinned to a fixed pair of physical cores, and a consumer-grade SD card bind-mounted as the spill target.
Together these constraints make the container \emph{resemble}, not replicate, a memory-constrained edge device: they reproduce its resource limits, not its hardware, and Section~\ref{sec:threats} states what this does and does not license.
Table~\ref{tab:edge-profile} summarizes the resulting profile.

\begin{table}[ht]
\centering
\small
\begin{tabular}{|>{\raggedright\arraybackslash}p{2.2cm}|>{\raggedright\arraybackslash}p{4.7cm}|>{\raggedright\arraybackslash}p{6.0cm}|}
\hline
\textbf{Resource} & \textbf{Container view} & \textbf{Mechanism} \\
\hline
Memory & 512 / 1024 / 2048\,MB, no swap & \texttt{--memory=Xm --memory-swap=Xm} \\
\hline
CPU & 2 cores, pinned & \texttt{--cpus=2 --cpuset-cpus=0,1,8,9} \\
\hline
Spill storage & SD card (SanDisk Extreme Pro, 128\,GB, ext4) & host \texttt{/mnt/sd-spill} bind-mounted at \texttt{/spill} \\
\hline
Input data & RAM-backed tmpfs & \texttt{/dev/shm} bind-mounted at \texttt{/input} (read-only) \\
\hline
\end{tabular}
\caption{Container device profile: the four imposed constraints that emulate the
resource envelope of a memory-constrained edge device, not a specific device.
The memory cap is varied per experiment (512, 1024, 2048\,MB); the other three
are held fixed.}
\label{tab:edge-profile}
\end{table}

Input data sits on a RAM-backed filesystem, keeping input reads off the SD card.
The spiller owns all SD card traffic: writes during eviction, reads during restore before emission.
Cells run sequentially: the SD card is a single physical device and concurrent spill I/O across cells would contaminate the per-cell latency measurements.
Each cell writes one row with 34 measured columns.

% [6.6]
Correctness is verified with output checksums against the in-memory baseline.
The gated binary itself is verified with a levers-off regression arm, which must behave identically to the unmodified engine.

The backend, the levers, and the instrumentation are in place.
The next chapter evaluates them.
